\documentclass[10pt,a4paper]{amsart}
\usepackage[margin=19mm]{geometry}
\usepackage{amsmath,amssymb,mathtools}
\usepackage[scr=rsfs,cal=euler]{mathalfa}
\usepackage{xfrac}
\usepackage[unicode,hidelinks,bookmarks=false]{hyperref}
\newcommand{\ie}{i.e.\ }
\newcommand{\cf}{cf.\ }
\newcommand{\resp}{respectively\ }
\newcommand{\Subsection}{\subsection}
\providecommand{\nobreakdash}{\nobreak\hbox{-}}
\setlength{\emergencystretch}{2em}
\multlinegap0pt
\usepackage[shortlabels]{enumitem}
\usepackage{mathtools}
\usepackage{amssymb}
\usepackage{xcolor}
\newtheorem{lem}{Lemma}
\newtheorem{prop}{Proposition}
\newcommand{\R}{\mathbb{R}}
\newcommand{\diam}{\operatorname{diam}}
\newcommand{\osc}{\operatorname{osc}}
\title{Lipschitz solvability of prescribed Jacobian and divergence for singular measures}
\author{Luigi De Masi, Andrea Marchese}
\date{Source: arXiv:2603.28912v1 (CC BY 4.0)}
\begin{document}
\maketitle

\noindent\emph{Source and licence.} Lipschitz solvability of prescribed Jacobian and divergence for singular measures. Version arXiv:2603.28912v1 (CC BY 4.0), distributed by arXiv under the Creative Commons Attribution 4.0 International licence, http://creativecommons.org/licenses/by/4.0/, as recorded in the arXiv licence metadata for this exact version and shown on the abstract page https://arxiv.org/abs/2603.28912v1. Full licence notice: \texttt{licenses/arxiv-2603.28912-CC-BY-4.0.txt}.\par\smallskip

\noindent\textit{Editorial scope.} Counted unit: the article's prop prop:local (source0.tex lines 348--364). The article's complete original proof of it is delivered with it (source0.tex lines 366--546, one \texttt{proof} environment of 181 lines). No mathematics is written, repaired or re-derived: the statement and the proof are the article's own lines, in the article's own notation, and no retained line is substituted or re-flowed.  Retained uncounted as the source-local context the proof needs: lines 313--322.  Nothing was omitted from any retained span, so the package carries no omission record.  No retained line contains a citation, so the wrapper carries no bibliography and no bibliography entry.  No retained line contains a figure environment, an image-inclusion command or drawing code, so no image is delivered and the package declares no asset dependency.  Editorial formatting only, in addition: an AMS \texttt{amsart} preamble that restates the article's own declarations and macros used by the retained lines, namely the environments \texttt{lem}, \texttt{prop} on separate counters, together with the standard packages those lines need.  The retained environments print in this document as Lemma 1, Proposition 1 in order of appearance, whereas the article prints the same items with the numbering of its own full document.  The complete native source of the article as distributed in the arXiv e-print for this exact version is bundled unchanged as \texttt{sources/197477/source0.tex}.
\begin{lem}[Width function]\label{lem:width}
Let $e\in \mathbb{S}^{d-1}$, $\alpha\in(0,\pi/2)$ and let $E\subset\R^d$ be compact and $C(e,\alpha)$-null.
Then for every $\zeta>0$ there exists a function $\varphi\in C^\infty(\R^d)$ such that:
\begin{enumerate}[label=\rm(\roman*),leftmargin=2em]
\item $0\le \varphi(x)\le \zeta$ for every $x \in \R^d$;
\item $0\le \partial_e\varphi(x)\le 1$ for every $x \in \R^d$ and $\partial_e\varphi(x)=1$ for all $x \in E$;
\item $|\partial_v\varphi(x)| \le \dfrac{1}{\tan\alpha}$ for every $v \in e^\perp$ with $|v|=1$ and every $x \in \R^d$.
% \item $|d_{Z_e}\varphi(x)| \le \dfrac{1}{\tan\alpha}$ for every $x \in \R^d$.
\end{enumerate}
\end{lem}

\begin{prop}\label{prop:local}
Let $U\subset\R^d$ be open, let $\nu$ be a finite Radon measure on $U$, let $e\in \mathbb{S}^{d-1}$, let $\alpha\in(0,\pi/2)$,
and let $\lambda\in C(U)\cap L^\infty(U, \nu)$.
Assume that there exists a $C(e,\alpha)$-null set $F\subset U$ such that $\nu(U\setminus F)=0$.
Then for every $\eta>0$ and every $\delta>0$ there exist a compact set $K\subset U$ and a function $u\in C^1_c(U)$ such that
$$
\nu(U\setminus K)<\eta,
\qquad
\partial_e u(x)=\lambda(x)\quad\forall x\in K,
$$
and
$$
\|Du\|_{C^0(U)}\le (1+\delta)c_\alpha \|\lambda\|_{L^\infty(U,\nu)},
\qquad
\|u\|_{C^0(U)}\le \eta.
$$
\end{prop}

\begin{proof}
 Set $M \coloneqq \|\lambda\|_{L^\infty(U,\nu)}$.
If $M=0$, taking $u\equiv 0$ and any compact set $K\subset U$ with $\nu(U\setminus K)<\eta$ gives the conclusion.
Hence, we may assume $M>0$.
Fix $t\in \left(0,\frac{1}{2}\right)$ such that
\begin{equation}\label{eq:t_geom_piccola}
    \sum_{n=0}^{+\infty}t^n + \sum_{n=0}^{+\infty}t^{2n+4} =\frac{1}{1-t} + \frac{t^4}{1-t^2}< 1+\min\{\delta,\eta\}.
\end{equation}

By inner regularity, choose a compact set $K_0\subset F$ such that
$$
\nu(U\setminus K_0)<\frac{\eta}{2}.
$$
Since $K_0\Subset U$, there exists an open set $W\subset U$ with
$$
K_0\Subset W\Subset U.
$$
We will construct inductively a decreasing sequence of compact subsets of $W$
$$
K_0\supset K_1\supset K_2\supset\cdots
$$
and, for every $n\geq 0$, a sequence of functions $u_{n+1}\in C^1_c(W)$
such that, setting
$$
% U_0\coloneqq 0, \qquad U_n\coloneqq \sum_{m=1}^{n}u_{m}\qquad
r_n\coloneqq \lambda-\partial_e \sum_{k=1}^n u_k,
$$
the following estimates hold for every $n\ge 0$:
\begin{align}
\nu(K_n\setminus K_{n+1}) &< \frac{\eta}{2^{n+2}}, \label{eq:loss}\\
\|r_n\|_{C^0(K_n)} &\le t^n M, \label{eq:residual}\\
\|Du_{n+1}\|_{C^0(U)} &\le t^n M\bigl(c_\alpha+t^{n+4}\bigr), \label{eq:gradstep}\\
\|u_{n+1}\|_{C^0(U)} &\le \tau\,t^{2n+4}M, \label{eq:supstep}
\end{align}
where
$$
\tau\coloneqq \min\left\{1,\frac{\delta}{M}, \frac{\eta}{M}\right\}.
$$
Since $r_0=\lambda$, then \eqref{eq:residual} is true at $n=0$.\\

Assume by induction that $K_n$ and $u_n$ are already defined and satisfy the estimates.
Since $r_n$ is continuous on the compact set $K_n$, it is uniformly continuous, namely there exists $\rho_n>0$ such that
$$
|r_n(x)-r_n(y)|\le t^{n+1}M
\qquad\forall x,y\in K_n,\ |x-y|<\rho_n.
$$
Choose a finite covering
$$
K_n\subset \bigcup_{i=1}^{N_n} B(x_{n,i},\rho_n/2).
$$
Define a Borel partition of $K_n$ by
$$
E_{n,1}\coloneqq K_n\cap B(x_{n,1},\rho_n/2),
$$
and for $i\ge 2$,
$$
E_{n,i}\coloneqq \Bigl(K_n\cap B(x_{n,i},\rho_n/2)\Bigr)\setminus \bigcup_{m<i}E_{n,m}.
$$
Then the sets $E_{n,i}$ are pairwise disjoint, their union is $K_n$ and $\diam(E_{n,i})\leq \rho_n$ for every $i$.
Hence the oscillation of $r_n$ on $E_{n,i}$ satisfies
$$
\osc_{E_{n,i}}(r_n)\le t^{n+1}M.
$$
By inner regularity, for every $i$ choose a compact set $L_{n,i}\subset E_{n,i}$ such that $\nu(E_{n,i}\setminus L_{n,i})<\frac{\eta}{2^{n+2}N_n}$ and set
$$
K_{n+1}\coloneqq \bigcup_{i=1}^{N_n}L_{n,i}.
$$
Then $K_{n+1}$ is compact, $K_{n+1}\subset K_n$, and
$$
\nu(K_n\setminus K_{n+1})
=\sum_{i=1}^{N_n}\nu(E_{n,i}\setminus L_{n,i})
< \frac{\eta}{2^{n+2}}.
$$
This proves \eqref{eq:loss}.

Since the compact sets $L_{n,i}$ are pairwise disjoint and contained in $W$, there exist pairwise disjoint open sets $O_{n,i}\subset W$ with $L_{n,i}\Subset O_{n,i}\Subset W$.
Choose cutoff functions $\chi_{n,i}\in C^\infty_c(O_{n,i})$ with $\chi_{n,i}\equiv 1$ in a neighborhood of $L_{n,i}$ and such that $0\le \chi_{n,i}\le 1$.

Since each $L_{n,i}\subset K_0\subset F$ and $F$ is $C(e,\alpha)$-null,
thus for every $i$ there exists $\varphi_{n,i}\in C^\infty(\R^d)$ satisfying the conclusions of Lemma~\ref{lem:width} with
$$
E=L_{n,i},
\qquad
\|\varphi_{n,i}\|_{C^0(\mathbb{R}^d)}
\leq
\zeta
=\dfrac{\tau\,t^{n+4}}{1+\max_i \|D\chi_{n,i}\|_{C^0(W)}} .
$$
Choosing points $y_{n,i}\in L_{n,i}$, we can set $a_{n,i}\coloneqq r_n(y_{n,i})$ and define
$$
w_{n,i}\coloneqq a_{n,i}\,\chi_{n,i}\,\varphi_{n,i},
\qquad
u_{n+1}\coloneqq \sum_{i=1}^{N_n}w_{n,i}.
$$
Since the supports of $\chi_{n,i}$ are pairwise disjoint, the supports of $w_{n,i}$ are pairwise disjoint as well.
Therefore $u_{n+1}\in C^1_c(W)$ and
$$
\|Du_{n+1}\|_{C^0(W)} = \max_i \|Dw_{n,i}\|_{C^0(W)},
\qquad
\|u_{n+1}\|_{C^0(W)} = \max_i \|w_{n,i}\|_{C^0(W)}.
$$

By Leibniz' rule,
$$
Dw_{n,i}=a_{n,i}\bigl(\chi_{n,i}D\varphi_{n,i}+\varphi_{n,i}D\chi_{n,i}\bigr),
$$
hence
$$
\|Dw_{n,i}\|_{C^0(W)}
\le \|r_n\|_{C^0(K_n)}\bigl(c_\alpha+\|\varphi_{n,i}\|_{C^0(W)} \|D\chi_{n,i}\|_{C^0(W)}\bigr)
\le t^{n}M\bigl(c_\alpha+t^{n+4}\bigr),
$$
which is \eqref{eq:gradstep}.
Similarly,
$$
\|w_{n,i}\|_{C^0(W)}
\le |a_{n,i}|\,\|\varphi_{n,i}\|_{C^0(W)}
\le t^{n}M\cdot \tau t^{n+4}
= \tau\,t^{2n+4}M,
$$
so \eqref{eq:supstep} follows.

It remains to control the new residual on $K_{n+1}$.
Fix $x\in L_{n,i}$.
Since $\chi_{n,i}\equiv 1$ on a neighborhood of $L_{n,i}$, we have $D\chi_{n,i}(x)=0$.
Since the supports are disjoint, all other functions $w_{n,m}$ with $m\neq i$ vanish in a neighborhood of $x$.
Hence
$$
\partial_e u_{n+1}(x)=\partial_e w_{n,i}(x)=a_{n,i}\,\partial_e\varphi_{n,i}(x)=a_{n,i}.
$$
Therefore
$$
r_{n+1}(x)=r_n(x)-\partial_e u_{n+1}(x)=r_n(x)-a_{n,i}.
$$
Since $x,y_{n,i}\in E_{n,i}$, we obtain
$$
|r_{n+1}(x)|
\le \osc_{E_{n,i}}(r_n)
\le t^{n+1}M.
$$
Thus $\|r_{n+1}\|_{C^0(K_{n+1})}\le t^{n+1}M$, \eqref{eq:residual} is proved and the induction is complete.

To conclude the proof of the proposition, observe that, by \eqref{eq:supstep}, it holds
$$
\sum_{n=0}^\infty \|u_{n+1}\|_{C^0(U)}
\le \tau M \sum_{n=0}^\infty t^{2n+4}
\le \frac{\eta}{8}
< \eta.
$$
and, by \eqref{eq:gradstep},
$$
\sum_{n=0}^\infty \|Du_{n+1}\|_{C^0(U)}
\le M\sum_{n=0}^\infty t^{n}\bigl(c_\alpha+t^{n+4}\bigr)
\overset{\eqref{eq:t_geom_piccola}}{<} (1+\delta)c_\alpha M.
$$
Thus the series $\sum_{n=0}^\infty u_{n+1}$ converges, in the $C^1$-norm, to a function $u \in C^1_c(W)$ with
$$
\|u\|_{C^0(U)} < \eta,
\qquad
\|Du\|_{C^0(U)}\le (1+\delta)c_\alpha M.
$$
Finally set $K\coloneqq \bigcap_{n=0}^\infty K_n$.
Since $K_n$ are decreasing compact sets, $K$ is compact.
Moreover, by \eqref{eq:loss},
$$
\nu(U\setminus K)
\le \nu(U\setminus K_0)+\sum_{n=0}^\infty \nu(K_n\setminus K_{n+1})
< \frac{\eta}{2}+\sum_{n=0}^\infty \frac{\eta}{2^{n+2}}
=\eta.
$$
If $x\in K$, then $x\in K_n$ for every $n$, and by \eqref{eq:residual}
$$
|r_n(x)|=\left|\lambda(x)-\partial_e \sum_{k=1}^{n}u_k(x)\right|\le t^{n}M\to 0.
$$
Since $\sum_{k=1}^{n}u_k\to u$ in $C^1$ as $n \to +\infty$, we conclude that
$$
\partial_e u(x)=\lambda(x)
\qquad\forall x\in K,
$$
completing the proof.
\end{proof}

\end{document}
